\section{Faza III — konfrontacja z agentem zewnętrznym}
\label{sec:phase3}

Po turnieju głównym fazy II przeprowadzono walidację zewnętrzną. Sprawdzono, czy przewaga wariantu \textit{21 głębokościowego} nie jest jedynie efektem porównań między zaproponowanymi wariantami i czy utrzymuje się także w starciu z agentem spoza eksperymentu. Punktem odniesienia był agent referencyjny z konkursu \textit{Hearthstone AI Competition}~\cite{hearthstone_competition,hearthstoneai_botdownloads_2020}.

Do porównania wybrano agenta Sebastiana Millera (\textit{2-step lookahead}), zwycięzcę edycji 2020. Agent Millera nie jest prostym agentem parametrycznym, lecz reprezentuje rodzinę metod przeszukujących sekwencje ruchów, co stawia go wyżej pod względem złożoności decyzyjnej.

\subsection{Protokół porównawczy}

Przeprowadzono dwa niezależne turnieje:
\begin{enumerate}
    \item wariant \textit{21 głębokościowy} (reprezentant fazy II) przeciwko agentowi Millera,
    \item agent bazowy 21-parametrowy z wagami publikowanymi przez autorów podejścia referencyjnego~\cite{sabberstone_coreai_results} przeciwko agentowi Millera.
\end{enumerate}

Protokół pozostał identyczny jak w poprzednich fazach: \(9\) par talii, \(1000\) gier na parę talii dla każdej kolejności graczy, co daje łącznie \(18\,000\) partii na agenta.

\subsection{Wyniki i interpretacja}

\begin{table}[ht]
\centering
\small
\begin{tabular}{llrrr}
\toprule
Agent badany & Przeciwnik & Wygrane & Partie & Odsetek zwycięstw [\%] \\
\midrule
21 głębokościowy & Sebastian Miller & 8525 & 18000 & 47,36 \\
Agent bazowy (wagi publikowane) & Sebastian Miller & 6249 & 18000 & 34,72 \\
\bottomrule
\end{tabular}
\caption{Wyniki fazy III: porównanie wariantu głębokościowego i agenta bazowego względem agenta Sebastiana Millera.}
\label{tab:phase3-external-benchmark}
\end{table}

Tabela~\ref{tab:phase3-external-benchmark} pokazuje, że Miller wygrywa oba porównania --- ale skala jego przewagi jest zupełnie inna. Przeciwko wariantowi \textit{21 głębokościowemu} wynosi ona zaledwie \(5{,}28\) p.p., natomiast przeciwko agentowi bazowemu aż \(30{,}56\) p.p.

Różnica ta dobrze pokazuje jakość zaproponowanej modyfikacji. Przejście od agenta bazowego do wariantu \textit{21 głębokościowego} podnosi skuteczność o \(12{,}64\) p.p. w starciu z tym samym, silnym przeciwnikiem. Wynika z tego, że uzyskany rezultat nie jest tylko efektem porównań wewnętrznych z fazy II, ale ma praktyczne przełożenie na rywalizację z zewnętrznym przeciwnikiem.

Z perspektywy pytań badawczych faza III wzmacnia odpowiedź na RQ3 --- komponent planowania sekwencji akcji jest najważniejszym źródłem poprawy jakości agenta. Jednocześnie wyniki pokazują granice obecnej pracy --- mimo wyraźnego postępu agent konkursowy pozostaje silniejszy, co wskazuje kierunek dalszego rozwoju.
